\paragraph{The matched shape-by-rule experiment separates transfer from refinement.}
The matched optimizer experiment comprises 256 runs on 16 previously studied target graphs. It retains byte-identical descriptor anchors for each graph, together with the previously selected initial radius and model-shot count. The resolution-aware policy combines radius-dependent sampling with unresolved stopping, crossing isotropic and shaped regions with Hoeffding and empirical-Bernstein rules without retuning on these targets. The experiment consumes 9,458,244 optimization shots and 2,097,152 independent output-validation shots. Previously recorded SPSA, COBYLA and original Bernstein runs provide comparators with the same anchors and cap. The original Bernstein rule is our own implementation, not a reproduction of a published variance-aware optimizer. The comparison is exploratory and post hoc because these targets had already been used to diagnose the original method.
\paragraph{Returned quality and incremental gain answer different questions.}
At depth 2, the resolution-aware shaped Bernstein policy has a paired refinement gain of 0.065 (95\% interval 0.000 to 0.196) percentage points and changes 2/32 outputs.
The original shaped Bernstein policy has a gain of 0.032508 percentage points. From unrounded paired graph records, the resolution-aware minus original gain difference is 0.032665 (95\% interval 0.000000 to 0.097995) percentage points, and the mean shot difference is -24284.38 (95\% interval -27360.00 to -20536.75). The number of distinct targets with a changed output is 1 for the resolution-aware policy and 1 for the original policy. These depth-specific paired comparisons distinguish changes in quality from avoided expenditure.
The resolution-aware policy's paired gain difference from SPSA is -0.13 (95\% interval -0.21 to -0.05) percentage points; positive values favor the resolution-aware shaped Bernstein policy.
The resolution-aware policy's paired gain difference from COBYLA is -0.02 (95\% interval -0.15 to 0.09) percentage points; positive values favor the resolution-aware shaped Bernstein policy.
At depth 3, the resolution-aware shaped Bernstein policy has a paired refinement gain of 0.032 (95\% interval 0.000 to 0.095) percentage points and changes 1/32 outputs.
The original shaped Bernstein policy has a gain of 0.108996 percentage points. From unrounded paired graph records, the resolution-aware minus original gain difference is -0.077215 (95\% interval -0.192421 to 0.000000) percentage points, and the mean shot difference is -26243.31 (95\% interval -29216.00 to -22534.62). The number of distinct targets with a changed output is 1 for the resolution-aware policy and 2 for the original policy. These depth-specific paired comparisons distinguish changes in quality from avoided expenditure.
The resolution-aware policy's paired gain difference from SPSA is -0.33 (95\% interval -0.47 to -0.20) percentage points; positive values favor the resolution-aware shaped Bernstein policy.
The resolution-aware policy's paired gain difference from COBYLA is -0.16 (95\% interval -0.37 to 0.00) percentage points; positive values favor the resolution-aware shaped Bernstein policy.
Figure~\ref{fig:refinement} reports returned quality, gain over each initial anchor, actual shots and the number of changed outputs. Figure~\ref{fig:efficiency} follows the output available after each completed measurement batch, without interpolating between separate budget-cap means. Initialization is available at zero shots, and an unfinished batch leaves the preceding output available. Target curves record first-ever attainment, retain failed runs in their denominator and count initial successes at zero cost. The paired difference between shaped policies uses pointwise graph-block bootstrap bands. Flat tails after termination show that the output remains available without further expenditure. Changing a parameter vector does not guarantee objective improvement for an unconstrained baseline. The reference is the best-found depth-matched objective, not a certified optimum; negative deficits are retained.
\begin{table}[t]
\caption{Resolution-aware refinement and matched original comparators. Values are graph-mean refinement gains and signed reference deficits in percentage points, followed by actual mean optimization shots in thousands. All policies use the same descriptor anchor for each graph. Each output also incurs 8,192 independent validation shots. RA denotes resolution-aware refinement; H and EB denote Hoeffding and empirical Bernstein. Original EB is our own comparator. The figures and text give paired 95\% bootstrap intervals obtained by resampling graphs within each fixed bank.}\label{tab:refinement}
\centering\small\setlength{\tabcolsep}{4pt}
\begin{tabular}{lrrrrrr}\toprule
 & \multicolumn{2}{c}{Gain} & \multicolumn{2}{c}{Deficit} & \multicolumn{2}{c}{Shots ($10^3$)} \\
Policy & $p=2$ & $p=3$ & $p=2$ & $p=3$ & $p=2$ & $p=3$ \\ \midrule
Initialization only & 0.000 & 0.000 & 0.407 & 0.768 & 0.000 & 0.000 \\
Original H, isotropic & 0.000 & 0.000 & 0.407 & 0.768 & 64.768 & 64.512 \\
Original EB, isotropic & 0.000 & 0.034 & 0.407 & 0.735 & 64.768 & 64.512 \\
Original EB, shaped & 0.033 & 0.109 & 0.375 & 0.659 & 64.768 & 64.512 \\
RA H, isotropic & 0.000 & 0.000 & 0.407 & 0.768 & 34.048 & 34.560 \\
RA H, shaped & 0.000 & 0.000 & 0.407 & 0.768 & 34.048 & 34.560 \\
RA EB, isotropic & 0.029 & 0.000 & 0.378 & 0.768 & 40.626 & 38.976 \\
RA EB, shaped & 0.065 & 0.032 & 0.342 & 0.736 & 40.484 & 38.269 \\
SPSA & 0.196 & 0.359 & 0.211 & 0.409 & 64.512 & 64.512 \\
COBYLA & 0.083 & 0.197 & 0.324 & 0.572 & 22.720 & 30.752 \\
\bottomrule\end{tabular}\end{table}
\newcommand{\ResolutionResults}{%
\paragraph{Stopping status distinguishes limited resolution from adverse evidence.}
The factorial records the following terminal statuses: \texttt{resolution\_limited}: 256. An unresolved comparison provides insufficient evidence to classify the direction. A resolution-limited stop records the end of the prescribed model or decision schedule without a resolved update. It does not establish stationarity, optimality, or the impossibility of a useful update under another allocation.
The stopping reasons are \texttt{acceptance\_look\_cap}: 230, \texttt{next\_acceptance\_look\_unaffordable}: 26.
}
\newcommand{\HistoricalDiagnostics}{%
\paragraph{Exact-data proposals isolate the affine model from sampling error.}
The offline diagnostic uses 2560 sampled fits across 8 fixed target graphs, both depths, both geometries, five radii and four shot counts. The design is held fixed within each target, depth and geometry. Of the 160 distinct exact-data fits, 84 improve the objective and 79 have a positive sufficient-decrease margin. At 1,024 shots per point, 292/640 sampled proposals improve and 233/640 have a positive sufficient-decrease margin. These counts describe dependent proposals on the chosen diagnostic grid; they are not counts of independent graphs. Figure~\ref{fig:components} separates sampling error around the exact affine fit from total error relative to the central-finite-difference reference gradient (step $10^{-5}$). The reference derivative is evaluated only offline. Neither lower fitting error nor better conditioning alone establishes useful refinement.
The model diagnostics require 83,558,400 simulated observations. The frozen-endpoint replay generates 2,684,354,560 simulated observations, with each stored stream shared across confidence rules and caps. These offline diagnostics use ideal sampling; their cost is reported separately from online optimization and is not counted as measurement savings.
At radius 0.025, 32/32 exact-data proposals have positive sufficient-decrease margin.
At radius 0.4, 0/32 exact-data proposals have positive sufficient-decrease margin.
\paragraph{Frozen endpoint replays measure the cost of decision resolution.}
The replay fixes 160 sampled-model proposals by target, depth, geometry and radius at 1,024 model shots per point, using the first diagnostic replicate. Each condition is replayed 128 times with independent endpoint streams and common error allocation across confidence rules. The extended replay uses a 131,072-shot endpoint-pair cap and five prescribed looks. Figure~\ref{fig:acceptance-power} compares all budget prefixes under both four-look and five-look allocations. Acceptance-stage cost includes unresolved runs at their caps. The analytical Hoeffding proposition gives an upper bound for its decision rule, not a universal QAOA shot lower bound.
For Hoeffding, the estimated mean acceptance probability over all fixed proposals is 0.000, and the mean unresolved fraction is 0.687.
Among the 58 positive-margin proposals, the estimated mean acceptance probability is 0.000, with mean endpoint cost 131,072 shots.
For Bernstein, the estimated mean acceptance probability over all fixed proposals is 0.104, and the mean unresolved fraction is 0.446.
Among the 58 positive-margin proposals, the estimated mean acceptance probability is 0.288, with mean endpoint cost 128,931 shots.
\paragraph{Paired regime comparisons delimit the evidence.}
Figure~\ref{fig:generality} retains negative effects in comparisons of shaped Bernstein with both conventional optimizers by vertex count and with SPSA by depth within graph families. The two reference banks belong to different target panels, so these comparisons do not identify a bank effect. All intervals average shot repetitions within graph and resample graphs within the two fixed banks (10,000 resamples, seed 5713); they are descriptive and unadjusted. This optimizer experiment provides no fresh holdout validation of the resolution-aware policy, large-graph scaling result, structural-shift test, or hardware-noise validation.
}
\newcommand{\RadiusResults}{%
\paragraph{Unresolved stopping often precedes increased model allocation.}
In the matched shape-by-rule experiment, 31/256 runs fit a model away from the initial radius (32/288 fits), and 28/256 runs allocate more than 256 shots per point (28/288 fits). Every resolution-aware Hoeffding run fits one initial model and completes one full four-look endpoint comparison before stopping with an unresolved decision. Its costs are therefore exactly $5(256)+2(16384)=34048$ at depth two and $7(256)+2(16384)=34560$ at depth three. The savings on these Hoeffding trajectories come from avoiding continuation; they do not demonstrate an improvement from changing the model-shot allocation.
\paragraph{The implemented schedule is tested against its variance prediction.}
The mechanism study uses 8 original diagnostic graphs from panel zero and 8 newly generated targets that are nonisomorphic to all previous bank, validation and test graphs. Figure~\ref{fig:components} displays shaped models on the original eight, averaging depths two and three equally within each graph before bootstrap resampling. The saved protocol records every graph identity and its split membership.
Central-difference gradients at steps $10^{-5}$ and $5\times10^{-6}$ differ by at most $1.39\times10^{-10}$ in Euclidean norm across all targets and depths. This agreement supports the numerical reference used to assess bias; the derivative remains a numerical approximation, rather than an analytic exact gradient.
The 1920 combinations of target, depth, shape, radius and shot count each use 128 independent sampled fits. Root-mean-square error is computed by averaging squared Euclidean errors, including across graphs, and then taking the square root; it is not the arithmetic mean of error norms. Exact-data bias floors and the propagated covariance predict total mean-square error (MSE). The implemented schedule uses 4096, 1024, 256, 64 and 16 shots per point at radii 0.025, 0.05, 0.1, 0.2 and 0.4. The total model cost includes all $2p+1$ fitting points. Graph-level bands are pointwise 95\% intervals conditional on the fixed bank and sampled designs; they do not provide simultaneous radius-wise coverage.
Across the fresh-graph model grid, the ratios of mean observed to mean predicted sampling and total MSE are 1.016122 and 1.014871, respectively. On the implemented schedule alone, the corresponding ratios are 1.009032 and 1.002424. Both summaries assign equal weight to fresh targets, depths and shapes. Agreement checks the standard conditional second-moment identity using exact simulator means, variances and numerical reference derivatives. This offline consistency check requires exact simulator inputs. Gaussian positive-margin prediction additionally approximates the distribution of fitted objective means and the nonlinear proposal map.
}
\newcommand{\PowerResults}{%
\paragraph{Online power uses its own four-look confidence allocation.}
Figure~\ref{fig:acceptance-power} uses fixed cohorts of 58 positive-margin proposals and the 102 nonpositive-margin proposals, with equal proposal weights at every budget. The online row uses four prescribed looks up to 16,384 shots per endpoint; the extended row uses five looks up to 65,536. The online acceptance panel is magnified to 0--0.06; the extended acceptance panel spans 0--1. Budgets are plotted on a base-two logarithmic axis. Both schedules are recomputed from the same stored independent streams, without additional sampled observations. Shading gives conservative pointwise 95\% Hoeffding bounds for each fixed-cohort replay mean, allowing different endpoint distributions across proposals. Those intervals quantify Monte Carlo replay uncertainty conditional on the fixed proposals, not uncertainty over target graphs.
At the 32,768-shot pair cap, the online four-look Hoeffding rule has a mean positive-margin acceptance frequency 0.0000 (pointwise interval 0.0000 to 0.0158), an unresolved fraction of 1.0000, and a mean actual endpoint expenditure of 32,768.00 shots.
At the 32,768-shot pair cap, the online four-look Bernstein rule has a mean positive-margin acceptance frequency 0.0233 (pointwise interval 0.0075 to 0.0391), an unresolved fraction of 0.9767, and a mean actual endpoint expenditure of 32,759.17 shots.
At the 131,072-shot pair cap, the extended five-look Hoeffding rule has a mean positive-margin acceptance frequency 0.0000 (pointwise interval 0.0000 to 0.0158), an unresolved fraction of 1.0000, and a mean actual endpoint expenditure of 131,072.00 shots.
At the 131,072-shot pair cap, the extended five-look Bernstein rule has a mean positive-margin acceptance frequency 0.2877 (pointwise interval 0.2720 to 0.3035), an unresolved fraction of 0.7123, and a mean actual endpoint expenditure of 128,931.31 shots.
For a single pair, observing zero acceptances in 128 independent replays gives an exact one-sided 95\% binomial upper limit of 0.02313 (2.313\%). Zero observed frequency therefore does not establish zero probability or empirically validate an analytical probability of order $10^{-6}$.
}
\newcommand{\MechanismResults}{%
\paragraph{A separate factorial isolates sampling allocation from unresolved stopping.}
The sampling-by-stopping experiment contains 256 shaped empirical-Bernstein runs on the same 16 targets. It crosses fixed and radius-dependent model-shot counts with stopping and continuing after unresolved evidence. All four cells share anchors, shape, model family, error allocation, shot cap, initial settings and paired seeds. Continuation retains the radius and refits with fresh data; the radius contracts only after certified rejection. This differs from the original policy, which also contracted after unresolved evidence. Figure~\ref{fig:attribution} separates gain and expenditure for these controlled policies.
\begin{table}[t]
\caption{Effects of model-shot allocation and stopping after unresolved evidence. Gains are percentage points over the common anchor; costs are actual optimization shots in thousands. Every row uses shaped empirical-Bernstein refinement at a 65,536-shot cap, with 8,192 independent validation shots per output. Two shot repetitions are averaged within each of 16 target graphs.}\label{tab:attribution}
\centering\small\begin{tabular}{lrrrr}\toprule
 & \multicolumn{2}{c}{Gain (pp)} & \multicolumn{2}{c}{Shots ($10^3$)} \\
Allocation / unresolved policy & $p=2$ & $p=3$ & $p=2$ & $p=3$ \\ \midrule
Fixed / continue & 0.0365 & 0.0723 & 64.7680 & 64.5120 \\
Fixed / stop & 0.0365 & 0.0723 & 39.0160 & 37.5840 \\
Radius / continue & 0.0365 & 0.0723 & 64.7752 & 63.8771 \\
Radius / stop & 0.0365 & 0.0723 & 39.9538 & 38.0696 \\
\bottomrule\end{tabular}\end{table}
All four cells return byte-identical parameter vectors in 64/64 matched graph--depth--seed cases. For this controlled sample, continuation and radius-dependent allocation leave returned quality unchanged, while stopping reduces expenditure.
At depth 2, with fixed model-shot counts, the paired difference in mean shots for stopping minus continuation is -25752.00 (95\% interval -28512.00 to -22440.00), a 39.76\% reduction in mean optimization shots (paired graph-bootstrap interval 34.65\% to 44.02\%).
At depth 2, with radius-dependent model-shot counts, the paired difference in mean shots for stopping minus continuation is -24821.44 (95\% interval -28181.44 to -20789.44), a 38.32\% reduction in mean optimization shots (paired graph-bootstrap interval 32.09\% to 43.73\%).
At depth 2, with stopping after unresolved evidence, radius-dependent minus fixed allocation gives a difference in mean gain of 0.00 (95\% interval 0.00 to 0.00) percentage points and a difference in mean shots of 937.81.
At depth 2, with continuation after unresolved evidence, radius-dependent minus fixed allocation gives a difference in mean gain of 0.00 (95\% interval 0.00 to 0.00) percentage points and a difference in mean shots of 7.25.
At depth 3, with fixed model-shot counts, the paired difference in mean shots for stopping minus continuation is -26928.00 (95\% interval -29952.00 to -23152.00), a 41.74\% reduction in mean optimization shots (paired graph-bootstrap interval 36.17\% to 46.43\%).
At depth 3, with radius-dependent model-shot counts, the paired difference in mean shots for stopping minus continuation is -25807.50 (95\% interval -29952.00 to -20926.00), a 40.40\% reduction in mean optimization shots (paired graph-bootstrap interval 32.94\% to 46.43\%).
At depth 3, with stopping after unresolved evidence, radius-dependent minus fixed allocation gives a difference in mean gain of 0.00 (95\% interval 0.00 to 0.00) percentage points and a difference in mean shots of 485.62.
At depth 3, with continuation after unresolved evidence, radius-dependent minus fixed allocation gives a difference in mean gain of 0.00 (95\% interval 0.00 to 0.00) percentage points and a difference in mean shots of -634.88.
Percentage savings are the difference in graph-mean costs divided by the continuation mean; both costs are resampled together within each fixed bank. These reductions concern the specified continuation policy, which retains the radius, at the shared cap. After a first complete test and two initial-radius models, only 30,208 shots at depth two or 29,184 at depth three remain for the next endpoint pair, less than the 32,768 shots required by its final look. The comparison does not establish how best to use the remaining budget.
\paragraph{Fresh-graph prediction is simulator-informed and conditional.}
All fitting maps, covariance predictions and Gaussian proposal predictions were frozen before the independent multinomial observations were generated. The fresh targets use the same fixed panel-zero donor bank, with no tuning on these targets. The fresh-graph study compares 256 Gaussian proposal draws per condition with 32 independent sampled proposals along the implemented radius schedule. The mean absolute difference in estimated positive-margin probability across fresh shaped conditions is 2.017 percentage points (graph-bootstrap interval 1.450 to 2.524). This discrepancy includes both Gaussian approximation error and finite-sample simulation noise. The predictor uses exact simulator objective means, variances and trial expectations. The experiment therefore tests the mechanism with exact target information; it does not provide a regime selector based only on measurements, an optimizer benchmark on fresh graphs, or hardware validation.
Figure~\ref{fig:attribution} reports signed observed-minus-Gaussian residuals separately from the absolute-error summary. Error bars resample target graphs. Gray regions give an exact conditional Monte Carlo reference: under equal Gaussian and sampled success probabilities within a graph, its 32-sample success count is hypergeometric conditional on the combined successes in 32+256 draws. Convolving the eight graph distributions gives a central 95\% reference for their equally weighted signed difference. 10/10 plotted depth/radius differences fall within this reference. This pointwise null range is neither a confidence interval for systematic approximation error nor uncertainty over a graph population. A zero-success condition gives a degenerate conditional reference, not evidence of zero underlying probability.
The sampling-by-stopping experiment consumes 13,201,786 optimization shots and 2,097,152 independent validation shots. The mechanism observations require 5,367,398,400 simulated finite-shot outcomes; the Gaussian diagnostic uses 40,960 draws and constructing predictions and references uses 43,872 exact simulator queries, plus 10,240 exact objective evaluations to diagnose sampled proposals. These offline diagnostic costs are not credited as optimizer savings.
}
\newcommand{\JointDecisionResults}{%
\paragraph{Simple prediction controls clarify the role of exact target information.}
The post-hoc radius-only and radius-plus-depth controls estimate condition means from the eight original diagnostic graphs, then evaluate the same shaped conditions on the eight fresh targets. Neither control uses fresh outcomes to fit its condition means. The comparison asks whether Gaussian predictions that use exact target information describe variation beyond these simpler controls. It does not supply a prediction rule based only on measurements.
The Gaussian predictor has a mean absolute frequency difference of 2.02 (95\% interval 1.45 to 2.52) percentage points and a mean Bernoulli Brier score of 0.051326.
The radius-plus-depth predictor has a mean absolute frequency difference of 19.28 (95\% interval 12.70 to 26.62) percentage points and a mean Bernoulli Brier score of 0.139438.
The radius-only predictor has a mean absolute frequency difference of 20.12 (95\% interval 13.87 to 27.11) percentage points and a mean Bernoulli Brier score of 0.142427.
Under the within-condition equal-probability null, the exact conditional expectation of the absolute frequency difference is 2.1937 percentage points; 10,000 independent conditional null draws give a central 95\% reference of 1.6455--2.7783 percentage points for its mean. The observed mean of 2.0166 lies within this Monte Carlo reference. This result is consistent with finite simulation noise at the available resolution; it neither proves equality nor isolates Gaussian approximation error.
\paragraph{Matched proposal and decision measurements reveal a conditional operating window.}
Figure~\ref{fig:joint-utility} applies the four-look endpoint rules to the same fresh shaped conditions used for proposal prediction at each radius and depth, with the implemented 4096, 1024, 256, 64 and 16 model shots per point. Each condition retains all 256 Gaussian proposals with one independent endpoint replay each, and all 32 sampled proposals with eight independent endpoint replays each. These eight replays share one model fit, rather than representing eight new fits. Proposals of every margin sign and all unsuccessful decisions remain in the denominator. Accepted true gain is set to zero when a decision fails and retains its sign when accepted. Each attempt is charged for one model and its actual endpoint measurements. Graph-level intervals are pointwise and descriptive over the eight fresh targets conditional on the fixed bank; they do not treat endpoint replays as independent graphs.
At depth 2, radius 0.1 and the 32,768-shot pair cap, the sampled observations give a mean accepted true gain under Bernstein of 0.003877 percentage points (95\% graph interval 0.000000 to 0.011630), an acceptance frequency of 0.4395\%, and a mean total cost of 32,976.00 shots (95\% graph interval 32,104.00 to 33,672.00). Positive accepted gain occurs on 1/8 target graphs.
At depth 2, radius 0.1 and the 32,768-shot pair cap, the Gaussian forecasts give a mean accepted true gain under Bernstein of 0.005787 percentage points (95\% graph interval 0.000000 to 0.017362), an acceptance frequency of 0.6348\%, and a mean total cost of 33,096.00 shots (95\% graph interval 31,984.00 to 33,816.00). Positive accepted gain occurs on 1/8 target graphs.
At depth 3, radius 0.1 and the 32,768-shot pair cap, the sampled observations give a mean accepted true gain under Bernstein of 0.179885 percentage points (95\% graph interval 0.000000 to 0.471226), an acceptance frequency of 14.8438\%, and a mean total cost of 33,336.00 shots (95\% graph interval 31,584.00 to 34,528.00). Positive accepted gain occurs on 2/8 target graphs.
At depth 3, radius 0.1 and the 32,768-shot pair cap, the Gaussian forecasts give a mean accepted true gain under Bernstein of 0.178342 percentage points (95\% graph interval 0.000000 to 0.481870), an acceptance frequency of 14.4043\%, and a mean total cost of 33,408.00 shots (95\% graph interval 31,904.00 to 34,528.00). Positive accepted gain occurs on 2/8 target graphs.
At radius 0.025, the sampled Bernstein proposals have a mean accepted gain of 0.000000 percentage points at depth 2, with a mean total cost of 53,248 shots.
At radius 0.025, the sampled Bernstein proposals have a mean accepted gain of 0.000000 percentage points at depth 3, with a mean total cost of 61,440 shots.
The Hoeffding rule accepts no proposal in either simulated population at any tested allowance.
High positive-margin frequency alone therefore does not identify the largest certified gain: proposal magnitude, decision resolution and model expenditure jointly determine the outcome. This post-hoc diagnostic evaluates a single attempt conditional on exact target information and a fixed bank. The resulting gain does not measure complete-optimizer performance on held-out graphs or compare that performance with SPSA or COBYLA.
The joint diagnostic generates 1,342,177,280 new simulated endpoint outcomes across 40,960 replays and 23,056 exact distribution queries. Its model proposals are reused from the saved mechanism experiment; no new model measurements are generated. The per-attempt cost nevertheless charges the saved model allocation once. Rules and budget prefixes reuse each new endpoint stream, so their reported costs must not be summed as independent offline sampling expense.
}
